\documentclass[conference]{IEEEtran}
\usepackage[T1]{fontenc}
\usepackage{cite}
\usepackage{amsmath,amssymb}
\usepackage{graphicx}
\usepackage{booktabs,array}
\usepackage{xcolor}
\usepackage{url}
\usepackage[hidelinks]{hyperref}
\hypersetup{pdftitle={When Operator Accuracy Does Not Transfer to the Image},
pdfauthor={Sally Mohlakola},pdfsubject={Learned interior transport for faceted gemstones}}
\graphicspath{{figures/}}
\newcommand{\dg}{^{\circ}}
\newcommand{\E}{\mathbb{E}}
\newcommand{\R}{\mathbb{R}}
\newcommand{\spp}{\,\mathrm{spp}}
\pagestyle{plain}
\raggedbottom

\begin{document}
\title{When Operator Accuracy Does Not Transfer to the Image:\\
A Learned Interior Transport Operator for Faceted Gemstones}
\author{\IEEEauthorblockN{Sally Mohlakola}
\IEEEauthorblockA{Honours Research Report}}
\maketitle

\begin{abstract}
Neural appearance models can replace detailed light transport with a compact
learned representation. This study investigates whether an adaptation of that
approach reproduces the appearance of resolved, specular gemstones. A
direction-conditioned radiance distribution map is first evaluated and then
generalised to an interior transport operator conditioned on entry position,
direction, surface normal and wavelength. The operator predicts finite-depth
escape, an exit triangle, and an exit position and direction. Stochastic and
deterministic decoders are evaluated on three procedural diamond cuts using
spectral reference renders, independent render seeds, and a separation of
operator-mediated light from untouched surface reflections. Learned exit-triangle
prediction exceeds the tested lookup baselines, but this improvement does not
ensure image agreement. At matched rendering budgets, stochastic decoders achieve
pixel correlations of 0.136--0.293, while deep deterministic decoders achieve
0.038--0.168 despite improved validation triangle accuracy. Facet-averaged
measurements reveal that coarse appearance can nevertheless be recovered: on
the step cut, correlation reaches 0.899 while fine-scale agreement remains low.
Recorded ablations implicate both branch selection and within-branch placement;
they do not establish a universal limit on neural transport. The results support
a narrower conclusion: under the tested studio illumination, operator-level
accuracy and perceived sharpness are insufficient proxies for rendered fidelity.
Training convergence, illumination coverage and computational cost are discussed
with explicit separation of completed measurements from further validation.
\end{abstract}

\begin{IEEEkeywords}
neural rendering, spectral light transport, gemstones, mixture density networks,
image evaluation, Monte Carlo rendering
\end{IEEEkeywords}

\section{Introduction}
The appearance of a faceted diamond depends on where internally reflected and
refracted light emerges. Small changes in entry position or direction can
change the intersected facets and redirect a path. Narrow sources produce
high-contrast structure sensitive to these changes; dispersion adds wavelength
dependence. An approximation must therefore preserve both transmitted energy
and its placement.

Soh and Montazeri~\cite{soh2024} aggregate explicit cloth-fibre scattering into
a compact neural appearance model. This motivates an adaptation to gemstones,
where visible facets and their internal reflections are resolved in the image.
Averaging their angular structure can remove the features the renderer must
reproduce.

The research question is consequently: \emph{how faithfully can a compact
learned interior transport operator reproduce resolved gemstone appearance, and
which operator errors explain its image-space failures?} The study begins with
a radiance distribution map (RDM), diagnoses its missing spatial conditioning,
and evaluates an operator that returns the outcome of an entire interior path.
The aim is to test, rather than assume, whether predictive gains improve rendering.

The contributions are threefold. First, lookup and input-ablation experiments
quantify the benefit of spatial conditioning for exit-triangle prediction.
Second, a factorised operator supports stochastic and deterministic placement
while preserving explicit geometric exit constraints. Third, an evaluation
combining pixel statistics, geometric-facet averages, render-seed comparisons
and path-component separation exposes a mismatch between operator accuracy and
image agreement. Three cut families exhibit the pixel-level mismatch, while
the step cut provides an informative counterexample at a coarser spatial scale.

The conclusions concern the evaluated models, budgets and scene, rather than
all learned representations of specular media.

\section{Background and Related Work}
\textbf{Aggregated appearance.} The cloth model of Soh and
Montazeri~\cite{soh2024} learns an aggregate BSDF from traced fibre-level
transport and includes an importance-sampling scheme. The present study adapts
the idea of compressing internal scattering, not the cloth geometry or its
complete implementation. The initial RDM is a directional summary. The proposed
operator instead models a joint exit event conditioned on the entry state.
Consequently, its results should be understood as an investigation of this
adaptation, not a direct benchmark of the original cloth method.

\textbf{Transport through bounded objects.} Vicini et
al.~\cite{vicini2019} learn shape-adaptive subsurface scattering using a
conditional variational autoencoder and local shape descriptors. Their work
establishes a relevant precedent for predicting transport through a region.
Here the focus is a smooth dielectric with discrete reflection/refraction
choices and sharply directed exits. Classical gemstone rendering explicitly
models the relevant optical mechanisms~\cite{sun2000}; a spectral path tracer
provides the numerical reference in this study~\cite{pharr2023,wilkie2014}.

\textbf{Conditional distributions and representations.} Mixture density
networks~\cite{bishop1994} represent multiple possible outcomes for an input.
They are attractive when an entry state can escape by several internal routes.
However, representing distinct outcomes does not by itself guarantee accurate
branch probabilities or angular placement. Neural-network spectral
bias~\cite{rahaman2019} motivates investigating Fourier
features~\cite{tancik2020}; local multiresolution hash
encodings~\cite{muller2022ngp} provide another possible representation. The
reported Fourier tests are limited experiments, not a test of every frequency
encoding or network family.

\textbf{Preserving variation and evaluating images.} LEAN
mapping~\cite{olano2010} motivates retaining second moments rather than only
means. The preliminary RDM experiment borrows this statistical idea, without
reproducing LEAN's surface-filtering model. For image assessment, SSIM and
FLIP~\cite{wang2004,andersson2020} offer complementary structural and perceptual
criteria. The completed evaluation here instead emphasises linear-luminance
agreement, highlight placement and facet averages with explicit render-noise
baselines. No SSIM or FLIP measurements are claimed.

\section{From a Radiance Map to an Interior Operator}
\label{sec:rdm}
The initial RDM bins exiting light by direction for each incoming direction.
Its round-cut render lacks sharp internal structure. In 99.8\% of
well-populated bins, standard deviation exceeds the mean; the median ratio is
2.53. This identifies substantial within-bin variation, but not its placement.

A moment-based modification activates shading points with probability
$p=1/(1+r^2)$ at gain $1/p$, where $r$ is relative standard deviation.
Position-keyed activation keeps the pattern fixed. This preserves the intended
moments in expectation, but produces speckle rather than correctly located
reflections: the missing position--direction relationship is not recovered.

Direction-conditioned tables predict exit triangles at accuracy 0.068 with
32 cells and 0.098 with 2,048, against 0.047 for the most frequent triangle.
Conditioning on the entry triangle reaches 0.192. This motivates retaining
position and normal alongside direction and wavelength; it does not exclude
benefits from higher resolution in other representations.

\section{Method}
\label{sec:method}
\subsection{Scene, geometry and reference transport}
All principal comparisons use one studio scene: a perspective camera with a
$30\dg$ field of view, a grey diffuse ground plane, textured area lights and
two small intense sources. This is \emph{one lighting rig with several sources},
not a single lamp. The non-uniform illumination makes exit-direction errors
visible by assigning different radiance to different parts of the environment.
Under spatially and directionally uniform incident radiance, many lossless
transport differences would be much harder to observe.

Rendering uses the scalar spectral variant of Mitsuba
3~\cite{mitsuba3}. The reference samples Fresnel reflection or refraction,
applies wavelength-dependent refraction using a diamond Sellmeier model, and
traces each interior path until escape or a depth cap of 128. The estimator
samples the BSDF without next-event estimation. Absorption and body colour are
not modelled. The term \emph{analytic reference} identifies these explicit
transport rules; the finite-sample image is not asserted to be converged
physical ground truth.

\begin{table}[t]
\caption{Geometry and primary-view evaluation support. Geometric facets are
coplanar groups of mesh triangles. Retained facets cover at least 20 pixels.}
\label{tab:geometry}
\centering\small
\begin{tabular}{@{}lrrr@{}}
\toprule
Property & Round & Pear & Step \\
\midrule
Vertices & 34 & 33 & 58 \\
Triangles / operator classes & 64 & 62 & 112 \\
Geometric facets & 57 & 33 & 50 \\
Retained visible facets & 25 & 17 & 25 \\
Stone-mask pixels & 134,158 & 90,607 & 93,618 \\
Training records & 838,736 & 838,717 & 838,852 \\
Validation records & 209,656 & 209,712 & 209,696 \\
Invalid records excluded & 184 & 147 & 28 \\
\bottomrule
\end{tabular}
\end{table}

The three procedural meshes are a round brilliant, pear brilliant and step cut
(Table~\ref{tab:geometry}). The nominal round proportions use crown angle
$34.5\dg$, pavilion angle $40.75\dg$ and table fraction 0.56. Mesh validation
checks closedness, manifold edges, outward orientation, positive signed volume
and degenerate faces. The operator classifies \emph{triangles}, which are not
identical to optical facets: the 64-triangle round mesh contains 57 coplanar
facet groups. This distinction is retained in the evaluation. Geometry hashes
are stored with training data and checkpoints to detect unintended mesh
mismatches; deliberate transfer is a separate experiment.

\subsection{Conditional exit representation}
The entry vector $\mathbf{x}\in\R^{10}$ comprises position $\mathbf{p}$,
interior direction $\boldsymbol\omega$, entry normal $\mathbf{n}$ and wavelength
$\lambda$. Positions are divided by the mesh radius and wavelength is encoded
as $(\lambda-595)/235$ with $\lambda$ in nanometres. Let $e$ denote escape before
the depth cap, $f$ an exit triangle and $\mathbf{y}\in\R^4$ exit coordinates.
For escaping paths, the model factorisation is
\begin{equation}
 p(e=1,f,\mathbf{y}\mid\mathbf{x})=
 s_\theta(\mathbf{x})\,\pi_\theta(f\mid\mathbf{x})\,
 q_\theta(\mathbf{y}\mid\mathbf{x},f),
\end{equation}
with $p(e=0\mid\mathbf{x})=1-s_\theta(\mathbf{x})$. The triangle and coordinate
distributions are understood to be conditional on escape. This makes the
non-escape event explicit instead of assigning an artificial exit to a
truncated path.

Two coordinates encode barycentric logits. If triangle vertices are
$\mathbf{v}_{f1},\mathbf{v}_{f2},\mathbf{v}_{f3}$, then
\begin{equation}
 \boldsymbol\beta=\operatorname{softmax}(y_1,y_2,0),\qquad
 \mathbf{p}_{\rm out}=\sum_{j=1}^{3}\beta_j\mathbf{v}_{fj}.
\end{equation}
The remaining two coordinates are outgoing tangent slopes:
\begin{equation}
 \boldsymbol\omega_{\rm out}=
 \frac{\mathbf{n}_f+y_3\mathbf{t}_f+y_4\mathbf{b}_f}
 {\|\mathbf{n}_f+y_3\mathbf{t}_f+y_4\mathbf{b}_f\|}.
\end{equation}
These transformations place the exit on the selected triangle and in its
outward hemisphere. They do not ensure that the sampled event corresponds to a
physically valid internal route or that it cannot immediately hit another part
of the stone.

\subsection{Model variants and optimisation}
A two-layer, width-64 SiLU encoder feeds escape and triangle heads. A learned
16-dimensional triangle embedding conditions the exit-coordinate decoder. The
\emph{stochastic} model uses four diagonal Gaussian mixture components in
transformed coordinates~\cite{bishop1994}. The \emph{clone} predicts a single
coordinate vector with a smooth-L1 regression objective. The \emph{deep clone}
uses a four-layer coordinate decoder with a separate encoder. All variants still
sample the exit triangle; ``deterministic'' refers only to placement given that
triangle. The round-cut stochastic, clone and deep-clone models contain
17,637, 15,557 and 28,741 trainable parameters, respectively.

For escape label $e_i$, training combines binary escape loss with triangle and
coordinate terms for escaping records:
\begin{equation}
 \mathcal L_i=\operatorname{BCE}(s_\theta(\mathbf{x}_i),e_i)
 +e_i\left[-\log\pi_\theta(f_i\mid\mathbf{x}_i)
 +\ell_{\rm coord}(\mathbf{y}_i)\right].
\end{equation}
The stochastic coordinate term is mixture negative log-likelihood; the clone
term is smooth-L1 error. Their total objective values are therefore not directly
comparable. Adam~\cite{kingma2015} uses learning rate $10^{-3}$, batch size 256
and gradient clipping at norm 10. Stochastic runs use 80 epochs and deep-clone
runs 30; selection retains the lowest validation objective within each run.

These comparisons share the same training pool within a cut but change decoder
family, objective, and in the deep case capacity and encoder arrangement.
They are \emph{model comparisons}, not a single-variable causal ablation. A
separate same-checkpoint component-mean experiment isolates Gaussian sampling
spread more directly.

\subsection{Data and render integration}
Each principal pool contains 1,048,576 path records from 65,536 entry states.
Repeated outcomes capture alternative Fresnel decisions. Invalid records are
excluded, and all outcomes of an entry state are assigned together to training
or validation. Table~\ref{tab:geometry} gives the resulting record counts.
Validation selects checkpoints; separately recorded held-out tests are labelled
as such rather than conflated with these validation scores.

Camera-matched entries initially cover an azimuth range but not every object
orientation. Entry distribution is consequently a controlled part of the
experiment, not a claim of complete surface coverage. Section~\ref{sec:coverage}
examines the resulting pavilion blind spot.

During rendering, entry reflection or transmission is sampled using the
analytic Fresnel rule. Transmitted entries are collected into batches of 512.
Each sampled operator exit starts a new scene ray; non-escaping events
terminate, and exits that immediately re-enter the stone are discarded. The
outside scene and first-surface reflection use the same transport code in both
modes. No importance correction makes the learned replacement an unbiased
estimator of the original interior transport. It approximates that transport,
and invalid-exit rejection can add further bias.

\section{Evaluation Methodology}
\label{sec:evaluation}
\subsection{Sampling, masking and empirical noise baselines}
Principal renders have resolution $512\times512$ and a box filter. Four
independent seeds at $32\spp$ form a $128\spp$ average for each mode. Splitting
the seeds into two groups yields independent $64\spp$ estimates. The main
pixel comparison and reference self-comparison use these matching half-budgets;
displayed images and facet means use the full $128\spp$ average. This distinction
prevents a full-budget image score from being compared uncritically with a
half-budget reference score.

Metrics use a fixed mask $M$ of pixels whose central primary ray hits the stone.
Linear RGB is converted to luminance with coefficients $(0.2126,0.7152,0.0722)$.
No per-image brightness normalisation is applied before the reported energy
measurements. Background and ground-plane caustics lie outside the mask. Thus,
these results assess the visible stone, not the complete scene's transport.

Reference split-half agreement estimates sampling repeatability at that budget.
It is not a strict upper bound on another method's correlation and does not
establish convergence of the reference. Additional reference samples remain
necessary to distinguish a stable fine-scale error from an unresolved rare
highlight. Independent render seeds also do not replace independent training
seeds~\cite{pharr2023}.

\subsection{Complementary image measurements}
For neural luminance $N_p$ and reference luminance $A_p$, the reported measures
include Pearson correlation $\rho(N,A)$, signed relative energy error
\begin{equation}
 \epsilon_E=\frac{\sum_{p\in M}N_p-\sum_{p\in M}A_p}
 {\sum_{p\in M}A_p},
\end{equation}
and relative root-mean-square error
\begin{equation}
 \epsilon_R=\frac{\sqrt{|M|^{-1}\sum_{p\in M}(N_p-A_p)^2}}
 {\overline A_M}.
\end{equation}
Highlight overlap is $|H_N\cap H_A|/|H_A|$, where each set contains the
brightest 1\% of masked pixels. Contrast is the luminance standard deviation
divided by its mean. Correlation measures spatial association but can conceal
brightness scaling. Energy error measures integrated brightness but not
placement. Highlight overlap directly tests a small, visually important part of
the image and is correspondingly sensitive to sampling noise. The metrics are
interpreted jointly rather than combined into an arbitrary overall score.

\subsection{Facet averages and component separation}
\label{sec:facets}
To examine coarser structure, coplanar mesh triangles are grouped by normal and
plane offset. Each visible group covering at least 20 pixels contributes its
mean luminance; correlation and relative RMSE are computed across these means,
giving each retained facet equal weight. This spatial averaging reduces the
influence of individual noisy pixels but also removes within-facet structure.
It complements, rather than replaces, the pixel evaluation. Its equally
weighted brightness statistics should not be interpreted as total image energy.

Each camera path is additionally tagged according to whether it passes through
the interior replacement. The frame decomposes additively as
$I=I_{\rm operator}+I_{\rm untouched}$. The reference uses the corresponding
interior-path tag. This separates errors in approximated transport from direct
reflections and scene contributions computed by unchanged code. High agreement
in the untouched component is a useful control, but does not by itself validate
the learned branch's integration.

\subsection{Operator diagnostics and provenance}
Operator evaluation measures exit-triangle top-1 accuracy and angular error of
the decoded direction when supplied with the recorded true triangle. The latter
isolates placement from triangle selection. Within-triangle \emph{spread}
around a prediction's own mean is a different quantity: a zero-spread prediction
can still point in the wrong direction. Repeated transport outcomes also mean
that triangle accuracy is not a complete assessment of distribution quality.

Tables and plots are assembled from saved evaluation JSON files and checkpoint
validation histories. Historical diagnostics are identified separately from
the matched three-cut render set. The companion source package records source
paths and hashes. Training-loss curves and controlled timing experiments are
not reconstructed from validation-only logs or undocumented measurements.

\section{Results}
\label{sec:results}
\subsection{Spatial conditioning and learned prediction}
\begin{table}[t]
\caption{Recorded round-cut exit-triangle baselines and operator tests.
Nearest-neighbour results use 4,096 queries over 1,046,359 stored paths;
these are historical diagnostics, separate from Table~\ref{tab:validation}.}
\label{tab:baselines}
\centering\small
\begin{tabular}{@{}lr@{}}
\toprule
Predictor & Top-1 accuracy \\
\midrule
Most frequent triangle & 0.047 \\
Entry-direction table (32 cells) & 0.068 \\
Entry-direction table (2,048 cells) & 0.098 \\
Entry-triangle table & 0.192 \\
Nearest neighbour & 0.287 \\
Majority among eight neighbours & 0.343 \\
Stochastic operator & 0.380 \\
Clone operator & 0.410 \\
\bottomrule
\end{tabular}
\end{table}

Table~\ref{tab:baselines} shows improvement over the tested lookup baselines.
The clone's 0.410 accuracy is approximately 43\% above the nearest-neighbour
score and 20\% above the stronger eight-neighbour majority score. A lookup is
a useful non-parametric comparator, not an upper bound on learnability.
Differences in evaluation sample support should also be retained when
interpreting rounded historical scores.

In the recorded 60-epoch conditioning ablation, all arms use seed 23 and the
same entry-grouped split. Removing position reduces top-1 accuracy from 0.404
to 0.273; keeping only direction yields 0.117. Removing wavelength changes it
from 0.4042 to 0.4039. Position is therefore useful under this protocol, whereas
wavelength contributes little to this particular triangle-classification score.
The result does not demonstrate absence of dispersion in the physical system.

\subsection{Validation behaviour}
\begin{figure*}[t]
\centering
\includegraphics[width=0.98\textwidth]{validation_curves.pdf}
\caption{Saved validation histories for the three cuts. Solid curves use the
stochastic decoder; dashed curves use the deep deterministic decoder. Triangle
accuracy and triangle cross-entropy are comparable across decoder families;
total losses are not, because their coordinate objectives differ. These are
single-training-seed histories, not mean curves with uncertainty bands.}
\label{fig:validation}
\end{figure*}

\begin{table}[t]
\caption{Checkpoint-selection validation results from saved metrics.
S: stochastic; D: deep deterministic. Epoch is the selected checkpoint.}
\label{tab:validation}
\centering\small
\begin{tabular}{@{}llrrr@{}}
\toprule
Cut & Model & Epoch & Triangle accuracy & Triangle CE \\
\midrule
Round & S & 80 & 0.371 & 2.547 \\
      & D & 30 & 0.392 & 2.480 \\
Pear  & S & 73 & 0.311 & 2.821 \\
      & D & 29 & 0.338 & 2.736 \\
Step  & S & 80 & 0.305 & 3.156 \\
      & D & 30 & 0.328 & 3.067 \\
\bottomrule
\end{tabular}
\end{table}

Figure~\ref{fig:validation} and Table~\ref{tab:validation} show that learning
improves triangle prediction on all cuts. Deep models obtain higher selected
validation accuracy and lower triangle cross-entropy than their stochastic
counterparts. Several selected checkpoints occur at the final training epoch,
including both round and step stochastic runs and their deep variants. Their
curves therefore do not establish fully converged optimisation. This limits any
claim that insufficient training or capacity has been eliminated.

For the round stochastic model, predicted mean escape probability is 0.9984
against a measured validation escape fraction of 0.9975. This is close agreement
of \emph{rates}, not proof of correct classification of individual escape
events. Non-escape cases are rare, and an always-escape predictor already has
high accuracy. The recorded Brier score is about 0.00249; discrimination of the
rare class would require additional class-specific reporting.

\subsection{Sharper images have lower pixel agreement}
\begin{figure*}[t]
\centering
\includegraphics[width=\textwidth]{gia_three_heads.png}
\caption{Round-cut renders from the existing experiment: stochastic, clone,
deep clone and analytic reference at 128 spp. Deterministic placement restores
sharp features but does not reproduce their locations. Embedded panel scores
use all render seeds; Table~\ref{tab:images} uses 64-spp half-splits.}
\label{fig:round}
\end{figure*}

\begin{table*}[t]
\caption{Matched three-cut image evaluation. Each mode has four 32-spp seeds
(128 spp total); correlation, relative RMSE, highlight overlap and energy error
use 64-spp half-images. Contrast uses the full average. ``Reference halves''
compares independent reference estimates and is an empirical sampling baseline.}
\label{tab:images}
\centering\small
\begin{tabular}{@{}llrrrrr@{}}
\toprule
Cut & Decoder/comparison & Correlation $\uparrow$ & Rel. RMSE $\downarrow$ & Top-1\% overlap $\uparrow$ & Energy error & Contrast \\
\midrule
Round & Stochastic & 0.152 & 3.750 & 0.197 & $-0.8\%$ & 0.93 \\
 & Clone & 0.048 & 4.741 & 0.071 & $+26.4\%$ & 2.35 \\
 & Deep clone & 0.038 & 4.849 & 0.117 & $+13.8\%$ & 2.74 \\
 & Reference halves & 0.957 & 1.140 & 0.777 & $+0.6\%$ & 3.78 \\
\midrule
Pear & Stochastic & 0.136 & 1.834 & 0.007 & $+17.3\%$ & 0.89 \\
 & Deep clone & 0.054 & 3.795 & 0.002 & $+47.8\%$ & 2.31 \\
 & Reference halves & 0.899 & 0.684 & 0.595 & $+0.4\%$ & 1.48 \\
\midrule
Step & Stochastic & 0.293 & 2.353 & 0.059 & $+14.2\%$ & 0.84 \\
 & Deep clone & 0.168 & 2.956 & 0.057 & $+15.3\%$ & 1.77 \\
 & Reference halves & 0.942 & 0.817 & 0.650 & $-0.3\%$ & 2.37 \\
\bottomrule
\end{tabular}
\end{table*}

The main comparison is shown in Fig.~\ref{fig:round} and
Table~\ref{tab:images}. Under the tested rig, deterministic decoders produce
sharper, higher-contrast images but lower pixel correlation and higher relative
RMSE on every cut. Round deep-clone correlation is 0.038 versus 0.152 for the
stochastic model; the corresponding pear values are 0.054 and 0.136, and step
values are 0.168 and 0.293. The reference half-correlations are much higher,
between 0.899 and 0.957.

The stochastic round image nearly preserves integrated energy, yet its contrast
is only 0.93 against 3.78 for the reference. Energy agreement therefore does not
imply recovery of structure. Conversely, the pear deep clone exceeds reference
contrast (2.31 versus 1.48) and has a $+47.8\%$ energy error. A sharper-looking
image is not necessarily an improvement in either placement or radiometry.
Highlight overlap also remains well below reference self-overlap, independently
supporting the observation of misplaced bright structure.

\begin{figure}[t]
\centering
\includegraphics[width=\columnwidth]{pear_heads.png}\\[2pt]
\includegraphics[width=\columnwidth]{step_heads.png}
\caption{Pear (top) and step (bottom): stochastic and deep-clone results beside
their references. Broad shape and brightness organisation can survive even when
fine internal reflections are incorrect. Images use 128 spp.}
\label{fig:cuts}
\end{figure}

\subsection{Coarse agreement is a distinct result}
\begin{figure*}[t]
\centering
\includegraphics[width=0.98\textwidth]{agreement_scales.pdf}
\caption{Image agreement at two spatial scales. Left: pixel half-split
correlation. Right: correlation of full-seed facet means, with each retained
facet weighted equally. Black marks indicate reference split-half agreement at
the corresponding scale. The step deep clone improves coarse facet agreement
while worsening pixel agreement. The different estimands and sample budgets
must not be interpreted as a pure denoising comparison.}
\label{fig:scales}
\end{figure*}

Facet averaging changes the interpretation (Fig.~\ref{fig:scales}). Stochastic
facet correlations are 0.399, 0.645 and 0.880 for round, pear and step cuts;
deep-clone values are 0.038, 0.414 and 0.899. Reference facet self-correlations
exceed 0.9995. Thus the step deep clone reproduces the ordering and magnitude of
coarse facet brightness substantially better than its low pixel correlation
alone suggests. Its facet relative RMSE is 0.333, compared with 0.376 for the
stochastic model.

This is an important qualification of the central finding. Operator improvement
does not invariably worsen \emph{every} image statistic: the effect depends on
spatial scale and the chosen quantity. Coarse agreement cannot establish correct
sparkle or internal reflection placement, because a facet mean intentionally
discards that structure. The combination of strong facet agreement and weak
pixel agreement provides a more precise diagnosis than declaring the entire
image either successful or unsuccessful.

\subsection{Branch selection and within-branch placement}
Historical held-out diagnostics distinguish spread from error. Within a recorded
exit triangle, stochastic samples have a median angular spread of approximately
$26.0\dg$, whereas the corresponding traced branches are almost collimated in
the sampled experiment. A clone has no stochastic coordinate spread, but its
mean direction error given the true triangle is still about $24.7\dg$.
Determinism therefore removes one source of blur without establishing correct
placement.

In a same-checkpoint intervention, sampling the mixture component but using its
mean coordinates reduces measured spread from $26.0\dg$ to $8.1\dg$; the
recorded image correlation falls from 0.174 to 0.153. The residual spread is
consistent with retained discrete sampling. This intervention is more directly
controlled than comparing separately trained decoder families.

\begin{figure}[t]
\centering
\includegraphics[width=\columnwidth]{oracle_facet.png}
\caption{Historical 256-spp diagnostic: clone, clone supplied with the traced
exit triangle, and reference. Supplying the correct triangle restores contrast
but does not correct the learned exit coordinates.}
\label{fig:oracle}
\end{figure}

Supplying the true exit triangle to the clone produces contrast 3.73 against
3.74 for its reference, but correlation remains 0.041
(Fig.~\ref{fig:oracle}). This implicates residual placement error even when
triangle selection is bypassed. It does not establish that correcting triangle
selection is generally harmful.

\subsection{Data, capacity, encoding and coverage}
Recorded data-ladder accuracy rises from 0.048 at 256 entry states to 0.332 at
16,384 and 0.347 at 65,536. This shows diminishing gains over the tested range,
not saturation under arbitrary data growth. A wider stochastic model improves
validation objective by about 0.33 nats before overfitting; a deeper decoder
reduces recorded placement error from $24.5\dg$ to $21.6\dg$. Three tested
Fourier scales yield errors of $25.8\dg$--$26.7\dg$ against the $21.6\dg$
control. These interventions do not remove the observed failure. Unequal
optimisation progress, a limited capacity range and single training seeds
prevent them from excluding those factors completely.

Nearest-neighbour placement is about $8.9\dg$ \emph{on queries where its chosen
triangle is correct}. This subset differs from evaluating a learned decoder
with the true triangle supplied for every query. The value is consequently a
diagnostic of locally recoverable exits, not a directly comparable universal
placement bound. Likewise, differing exits among nearby entries may reflect
both branch-boundary sensitivity and Fresnel branching; neighbourhood diversity
alone cannot prove an unsampleable discontinuity density.

\label{sec:coverage}
Camera-matched training initially supplies no pavilion entries. In a recorded
60-frame pear tumble, pavilion-facing poses reduce predicted escape to 46\%.
Adding uniformly directed entries with a disc aim offset places approximately
53\% of entries on the pavilion and restores recorded escape to 100\% across
the sequence. At a pavilion-facing pose, correlation remains 0.018 against
reference repeatability of 0.951. Coverage repairs a specific failure mode but
does not by itself recover placement. The tumble is a coverage diagnostic,
not a substitute for a held-out lighting and viewpoint benchmark.

\subsection{Transfer and apparent agreement from untouched light}
The categorical head is mesh-specific. A separate experiment maps a pear-trained
operator to an unseen round variant by matching facet normals. This changes
geometry and correspondence simultaneously, so it tests a particular retargeting
heuristic rather than geometry-independent generalisation.

\begin{figure}[t]
\centering
\includegraphics[width=\columnwidth]{transfer_split.png}
\caption{Retargeting diagnostic from the recorded experiment. A direct lamp
reflection on the crown contributes apparent agreement although it bypasses
the learned operator. Component-separated evaluation exposes this effect.}
\label{fig:transfer}
\end{figure}

Whole-stone correlations of approximately 0.37 and 0.35 initially suggest
improved agreement. However, about 90\% of a dominant crown highlight is
untouched light. Removing an $81\times81$ patch around it reduces the recorded
correlation to $-0.001$, and the operator component correlates at only 0.004.
On the default round experiment, analogous patch removal leaves a recorded
correlation of 0.205 essentially unchanged. The lesson is methodological:
agreement produced by a shared analytic reflection is not evidence that the
learned interior transport transfers.

\subsection{Recorded computational cost}
\begin{table}[t]
\caption{Median recorded seconds per 512$\times$512, 32-spp render, over four
seeds. These historical runs were not a controlled hardware benchmark;
cross-run differences may include concurrent load.}
\label{tab:timing}
\centering\small
\begin{tabular}{@{}lrrr@{}}
\toprule
Mode & Round & Pear & Step \\
\midrule
Analytic & 657 & 485 & 497 \\
Stochastic & 1,269 & 955 & 746 \\
Deep clone & 767 & 725 & 655 \\
\bottomrule
\end{tabular}
\end{table}

Table~\ref{tab:timing} reports the median of the saved per-seed durations, not
the time to obtain the complete four-seed average. None of the recorded neural
configurations is faster than its cut's reference in these logs. Their
interpretation is limited by uncontrolled system load and incomplete hardware
provenance. Earlier project notes report 5--9~$\mu$s per batched query, but that
figure is insufficient for an end-to-end speed claim: input preparation,
sampling, decoding and scene-ray bookkeeping also contribute. No controlled
training-time, memory or throughput benchmark is asserted here.

\section{Discussion}
\subsection{Why predictive gains need not improve the image}
The image depends on the exit distribution weighted by the radiance encountered
after escape. Schematically, for entry state $\mathbf{x}$ and exit state $z$,
\begin{equation}
 L(\mathbf{x})=\int L_{\rm ext}(z)\,p(z\mid\mathbf{x})\,\mathrm{d}z.
\end{equation}
A training objective that weights all recorded exits similarly need not
prioritise the exits that connect to narrow bright sources. A small angular
error near a source boundary can have a large image effect, whereas a larger
error between equally dark directions may contribute little. Triangle accuracy
and average angular error are therefore useful but incomplete proxies for
rendered fidelity.

A broad predicted distribution can spread energy over both correct and
incorrect regions, preserving some correlation while suppressing contrast. A
sharper but misplaced distribution can lose this overlap. The results are
consistent with this explanation, but do not establish a universal monotonic
relationship between blur and correlation. The facet-scale step result shows
explicitly that an improvement can be visible in one image statistic while
remaining absent in another.

The representation is one plausible contributor. Smooth conditional densities
approximate transport that contains sharply concentrated branches, and changes
of internal route may create rapidly varying exit behaviour. Spectral bias
provides a relevant motivation~\cite{rahaman2019}, but incomplete convergence
and the tested objectives remain alternative or interacting explanations.
Local encodings, route-aware representations and radiance-weighted objectives
are hypotheses for further work, not demonstrated solutions.

\subsection{What the current evidence supports}
The evidence supports three bounded conclusions. Spatial conditioning improves
exit-triangle prediction in the tested setting. Improved validation prediction
does not guarantee improved fine-scale image agreement. Finally, reporting
pixel, facet and operator-only measurements together yields a more informative
assessment than a single whole-image score. The study provides an empirical
failure analysis of a compact learned replacement, rather than a proof that
neural transport cannot represent specular gemstones.

\section{Limitations and Further Validation}
\label{sec:limits}
The study uses one rig, mostly one pose, simplified meshes and mesh-specific
categorical heads. Each principal configuration has one training seed; several
validation curves are still improving. Reference convergence remains unverified,
stone masking excludes ground caustics, and saved logs lack paired training-loss
curves. Historical timings are not controlled performance measurements.

Four proposed tests address these gaps: (1) return traced exits through the
learned operator's interface to measure integration error and rejection;
(2) repeat training with at least three seeds and adequate optimisation;
(3) test reference stability at 128, 512 and 2,048 spp, extending if necessary;
and (4) evaluate fixed models under broad and narrow-source lighting at several
views inside training coverage. These are future experiments, not reported
results.

For a more noise-tolerant supplementary evaluation, normalised mean absolute
error can be reported as
\begin{equation}
 \operatorname{NMAE}(N,A)=
 \frac{\sum_{p\in M}|N_p-A_p|}{\sum_{p\in M}A_p+\epsilon},
\end{equation}
where $\epsilon$ is a documented small denominator safeguard. Evaluating it
after matched $2\times2$, $4\times4$ and $8\times8$ block averaging would
separate spatial scales; native-resolution scores must remain visible so that
averaging does not hide the failure being studied. These NMAE and multiscale
scores have not yet been measured. Their uncertainty should be assessed with
independent render seeds, alongside reference self-comparisons, rather than
assuming any metric is immune to Monte Carlo noise.

A controlled performance study should record hardware, software, threads,
precision, warm-up, batching and model-loading policy. It should separate
network latency, complete operator-call latency and render time, and compare
image error at equal elapsed time. Mesh and checkpoint hashes, scene parameters
and exposure settings are also needed for reproducibility.

\section{Conclusion}
A direction-conditioned radiance map was generalised to a position-conditioned
interior transport operator for three procedural gemstone cuts. The learned
models improve exit-triangle prediction over the tested lookup baselines, but
the deterministic variants' higher validation accuracy does not translate into
better fine-scale image agreement under the evaluated studio rig. Highlight
overlap and relative RMSE support the correlation-based finding, while
facet-averaged evaluation reveals substantial coarse agreement on the step cut.
True-triangle and component-separation diagnostics further distinguish
placement error from branch selection and from apparent agreement due to
untouched reflections. The resulting contribution is a measured account of
where this approximation succeeds and fails. Stronger claims about
representation limits, generalisation or computational advantage require the
additional convergence, repeatability and performance controls identified here.

\begin{thebibliography}{99}
\bibitem{soh2024} G. Y. Soh and Z. Montazeri, ``Neural appearance model for cloth
rendering,'' \emph{Computer Graphics Forum}, vol. 43, no. 4, 2024,
doi: 10.1111/cgf.15156.
\bibitem{vicini2019} D. Vicini, V. Koltun, and W. Jakob, ``A learned shape-adaptive
subsurface scattering model,'' \emph{ACM Trans. Graph.}, vol. 38, no. 4,
Art. 126, 2019, doi: 10.1145/3306346.3322974.
\bibitem{sun2000} Y. Sun, F. D. Fracchia, and M. S. Drew, ``Rendering diamonds,''
in \emph{Proc. 11th Western Computer Graphics Symposium}, 2000.
\bibitem{pharr2023} M. Pharr, W. Jakob, and G. Humphreys,
\emph{Physically Based Rendering: From Theory to Implementation}, 4th ed.
MIT Press, 2023. [Online]. Available: \url{https://pbr-book.org/4ed/}
\bibitem{wilkie2014} A. Wilkie, S. Nawaz, M. Droske, A. Weidlich, and J. Hanika,
``Hero wavelength spectral sampling,'' \emph{Computer Graphics Forum}, vol. 33,
no. 4, pp. 123--131, 2014, doi: 10.1111/cgf.12419.
\bibitem{bishop1994} C. M. Bishop, ``Mixture density networks,'' Aston University,
Tech. Rep. NCRG/94/004, 1994.
\bibitem{rahaman2019} N. Rahaman \emph{et al.}, ``On the spectral bias of neural
networks,'' in \emph{Proc. 36th Int. Conf. Machine Learning}, PMLR, vol. 97,
pp. 5301--5310, 2019.
\bibitem{tancik2020} M. Tancik \emph{et al.}, ``Fourier features let networks learn
high frequency functions in low dimensional domains,'' in \emph{Advances in
Neural Information Processing Systems}, vol. 33, pp. 7537--7547, 2020.
\bibitem{muller2022ngp} T. M\"uller, A. Evans, C. Schied, and A. Keller,
``Instant neural graphics primitives with a multiresolution hash encoding,''
\emph{ACM Trans. Graph.}, vol. 41, no. 4, Art. 102, 2022,
doi: 10.1145/3528223.3530127.
\bibitem{olano2010} M. Olano and D. Baker, ``LEAN mapping,'' in \emph{Proc. ACM
SIGGRAPH Symp. Interactive 3D Graphics and Games}, pp. 181--188, 2010.
\bibitem{wang2004} Z. Wang, A. C. Bovik, H. R. Sheikh, and E. P. Simoncelli,
``Image quality assessment: From error visibility to structural similarity,''
\emph{IEEE Trans. Image Process.}, vol. 13, no. 4, pp. 600--612, 2004,
doi: 10.1109/TIP.2003.819861.
\bibitem{andersson2020} P. Andersson, J. Nilsson, T. Akenine-M\"oller,
M. Oskarsson, K. \AA str\"om, and M. D. Fairchild,
``FLIP: A difference evaluator for alternating images,''
\emph{Proc. ACM Comput. Graph. Interact. Tech.}, vol. 3, no. 2,
Art. 15, 2020, doi: 10.1145/3406183.
\bibitem{mitsuba3} W. Jakob \emph{et al.}, ``Mitsuba 3 renderer,'' 2022.
[Online]. Available: \url{https://mitsuba-renderer.org}
\bibitem{kingma2015} D. P. Kingma and J. Ba, ``Adam: A method for stochastic
optimization,'' in \emph{Proc. Int. Conf. Learning Representations}, 2015.
\end{thebibliography}
\end{document}
